	\begin{UseCase}{CU12.1.4}{Gestionar Puntos de extensión}{

		Una vez que el colaborador (\hyperlink{jefe}{Líder de Análisis} o \hyperlink{analista}{Analista}) ha registrado un \hyperlink{casoUso}{Caso de Uso}, podrá acceder a las acciones que corresponten a este, tal como la gestión de \hyperlink{entidadExtension}{Puntos de Extensión}. Cada caso de uso contiene la gestión de sus respectivos puntos de extensión. 
		
		Este caso de uso permite al colaborador visualizar en una tabla, los puntos de extensión que conforman un Caso de uso en específico, así como solicitar el registro de uno nuevo, modificar la información de algún Punto de extensión existente (teniendo la posibilidad de actualizar algún dato de registro), así como eliminarlo (en caso de que el punto de extensión por algún motivo ya no forme parte del Caso de Uso). \\
	}
	\UCitem{Actor}{\hyperlink{jefe}{Líder de análisis}, \hyperlink{analista}{Analista}}
	\UCitem{Propósito}{Proporcionar al actor un mecanismo para llevar el control de los puntos de extensión de un caso de uso.}
	\UCitem{Entradas}{Ninguna}
	\UCitem{Salidas}{\begin{itemize}
			\item \cdtRef{proyectoEntidad:claveProyecto}{Clave del proyecto}: Lo obtiene el sistema.
			\item \cdtRef{proyectoEntidad:nombreProyecto}{Nombre del proyecto}: Lo obtiene el sistema.
			\item \cdtRef{moduloEntidad:claveModulo}{Clave del Módulo}: Lo obtiene el sistema.
			\item \cdtRef{moduloEntidad:nombreModulo}{Nombre del Módulo}: Lo obtiene el sistema.
			\item \cdtRef{casoUso:claveCU}{Clave} del caso de uso: Lo obtiene el sistema. 
			\item \cdtRef{casoUso:numeroCU}{Número} del caso de uso: Lo obtiene el sistema. 
			\item \cdtRef{casoUso:nombreCU}{Nombre} del caso de uso: Lo obtiene el sistema.
			\item \cdtRef{entidadExtension}{Punto de extensión}: Tabla que muestra la \cdtRef{entidadExtension:observacionRevision}{causa}, la región de la trayectoria y el caso de uso que extiende de todos los registros.
			\item \cdtIdRef{MSG2}{No existe información}: Se muestra en la pantalla \IUref{IU6.1.4}{Gestionar Puntos de extensión} cuando no existen puntos de extensión registrados.
	\end{itemize}}
	\UCitem{Precondiciones}{\begin{itemize}
			\item Que exista al menos un caso de uso registrado.
			\item Que existe al menos un punto de extensión registrado.
	\end{itemize}}
	\UCitem{Postcondiciones}{Ninguna}
	\UCitem{Errores}{Ninguno}
	\UCitem{Tipo}{Secundario, extiende del caso de uso \UCref{CU12}{Gestionar Casos de uso}.}
\end{UseCase}
%--------------------------------------
\begin{UCtrayectoria}
	\UCpaso[\UCactor] Solicita gestionar los puntos de extensión presionando el botón \raisebox{-1mm}{\includegraphics[height=11pt]{images/Iconos/talt}} del caso de uso que desee de la pantalla \IUref{IU6}{Gestionar Casos de uso}
	\UCpaso[\UCsist] Obtiene la información de los puntos de extensión del caso de uso. \hyperlink{CU12-1-6:TAA}{[Trayectoria A]}
	\UCpaso[\UCsist] Ordena los puntos de extensión alfabéticamente con base en la causa de los mismos.
	\UCpaso[\UCsist] Muestra la información de los puntos de extensión en la pantalla \IUref{IU6.1.4}{Gestionar Puntos de extensión}. \label{CU12.1.6-P4}
	\UCpaso[\UCactor] Gestiona los puntos de extensión a través de los botones: \IUbutton{Registrar}, \editar, y \eliminar.
\end{UCtrayectoria}		
%--------------------------------------
\hypertarget{CU12-1-6:TAA}{\textbf{Trayectoria alternativa A}}\\
\noindent \textbf{Condición:} No existen registros de puntos de extensión.
\begin{enumerate}
	\UCpaso[\UCsist] Muestra el mensaje \cdtIdRef{MSG2}{No existe información} en la pantalla \IUref{IU6.1.4}{Gestionar Puntos de extensión} para indicar que no hay registros de puntos de extensión para mostrar. \label{CU12.1.6-TA1}
	\UCpaso[\UCactor] Gestiona los casos de uso a través del botón: \IUbutton{Registrar}. 
	\item[- -] - - {\em {Fin del caso de uso}}.%
\end{enumerate}
%--------------------------------------

\subsubsection{Puntos de extensión}

\UCExtenssionPoint{El actor requiere registrar un punto de extensión.}{Presionando el botón \IUbutton{Registrar} del paso \ref{CU12.1.6-P4} de la trayectoria principal o del paso \ref{CU12.1.6-TA1} de la Trayectoria alternativa A .}{\UCref{CU12.1.4.1}{Registrar Punto de extensión}}
\UCExtenssionPoint{El actor requiere modificar un punto de extensión.}{Presionando el icono \editar del paso \ref{CU12.1.6-P4} de la trayectoria principal.}{\UCref{CU12.1.4.2}{Modificar Punto de extensión}}
\UCExtenssionPoint{El actor requiere eliminar un punto de extensión.}{Presionando el icono \eliminar del paso \ref{CU12.1.6-P4} de la trayectoria principal.}{\UCref{CU12.1.4.3}{Eliminar Punto de extensión}}
